\chapter{绪论}

处理器并不是一个天然可信的黑盒。即使厂商公开了 ISA 手册、编译器后端和常见反汇编器，处理器真实接受的编码集合、工具链对该编码的解释、以及不同执行环境中的行为结果之间仍可能存在差异。这类差异既可能体现为狭义的 undocumented instruction，也可能体现为解码器误判、模拟器工件、特权相关行为或执行期异常；对系统软件验证、漏洞分析和第三方审计而言，它们都具有直接价值。

x86 上的 \tool{sandsifter} 工作已经证明，系统化的指令空间审计能够发现隐藏指令、硬件异常以及工具链缺陷 \citep{domas2017breakingx86}。RISC 方向上，\tool{iScanU} 进一步将这一路线推广到 RISC-V 与 ARMv8，并给出了 \ptracebackend\ 与 \memcage\ 两条用户态执行路径 \citep{Dofferhoff_2020}。与此同时，\tool{armshaker} 也说明，很多看似“隐藏指令”的现象实际上需要在 CPU、解码器、模拟器与系统软件之间做细致 triage \citep{Strupe_2020}。这些工作共同构成了本文的直接背景。

然而，现有方案在本文所关心的场景中仍有四个缺口。第一，多数原型以单 ISA 为主，缺少明确的多 ISA 统一工程骨架。第二，\memcage\ 与 \ptracebackend\ 往往被当作独立技巧，而不是同一系统中的互补后端。第三，容器或仿真结果与真实硬件证据之间常被分开叙述，导致可复现性边界不够清晰。第四，新架构 bring-up 仍然高度依赖手工改写后端，缺少“配置 + 生成 + 有界人工加固”的接入路径。

本文的核心思路是用统一编排框架把这些问题收束到一条主线上：前端以 Python 负责编排参数、worker 与日志聚合，中间层以共享执行核心管理“写指-执行-回收”的公共逻辑，后端以 \code{arch\_*.c} 封装架构差异；同时保留 \memcage\ 与 \ptracebackend\ 两种执行控制，并增加配置驱动的后端生成链路，把新 ISA 的初始接入从“完全重写”收缩为“规格输入 + 可运行基线 + 明确 TODO 清单”。

\placeholderfigure{fig:hero}{[图 1-1 待补]\\总体框架图\\左：写指-执行-信号-反汇编-分类闭环\\中：\memcage\ 与 \ptracebackend\ 的互补路径\\右：\code{arch-spec} 到 \code{arch\_*.c} 的生成链路\\下：Layer A/B/C 三层证据结构}{本文工具的总体结构与证据路径。该图将在终稿中替换为正式框架图，用于让读者在阅读方法细节前先理解统一编排、双后端和生成器之间的关系。}

围绕上述目标，本文试图回答三个问题。其一，如何在用户态环境中，对任意指令编码实现单条、可回收、可分类的执行测试。其二，如何在同一框架中组织 \memcage{} 与 \ptracebackend{}，并在多 ISA 场景下复用同一套编排逻辑。其三，当新增一种 ISA 时，能否将 bring-up 过程从完全手工重写后端收缩为配置驱动的基线生成问题。

本文的主要贡献如下。
\begin{enumerate}[leftmargin=2em]
  \item 设计并实现了一套面向 RISC-V 与 AArch64 的用户态隐藏指令分析框架，统一了写指、执行、信号捕获、反汇编与分类主循环。
  \item 在统一框架中组织了 \memcage\ 与 \ptracebackend\ 两类执行后端，并给出当前 AArch64 对齐子集上的 matched-subset 吞吐对比。
  \item 提出并实现了配置驱动的新架构生成链路，并以 MIPS64 作为首个案例验证“配置 + 生成 + 有界人工加固”的可行性。
  \item 采用 Layer A、Layer B、Layer C 的分层证据结构；当前稿已经完成 Layer A 与 Layer C，Layer B 保留为后续实机补证部分。
\end{enumerate}

当前最强结果包括三部分。第一，RISC-V 与 AArch64 的 Layer A 短程扫描已能稳定完成并落盘。第二，在 AArch64 的对齐子集上，\memcage\ 的 median 吞吐相对 \ptracebackend\ 高约 14.3\%。第三，MIPS64 的生成器案例已经从 36 行规格输入产出 497 行可编译后端，并完成 quick smoke。与此同时，本文并不把当前稿写成“已经证明所有主张”的定稿：真实硬件可复现性和 Layer A/Layer B 差异归因仍待后续实验补齐。

本文余下内容组织如下。第 2 章讨论相关工作与技术背景；第 3 章给出问题定义、分类口径与设计目标；第 4 章介绍系统设计；第 5 章说明实现细节；第 6 章报告实验与评估；第 7 章讨论局限；第 8 章总结全文并给出后续工作。
